\chapter{CORS --- Connecting Different Origins}

\section{Why the Browser Blocks the Request}

During local development the React app runs on
\texttt{http://localhost:5173} (Vite's default port) and the Express API runs
on \texttt{http://localhost:5000}. Even though both are ``localhost'', they
are two different \textbf{origins} --- an origin is the triple
\texttt{scheme + host + port}, and any difference in any one of those three
makes it a cross-origin request. Browsers enforce the
\textbf{Same-Origin Policy}: JavaScript running on one origin cannot read
responses from a different origin unless that origin explicitly says it is
allowed to. This is a browser-side security rule, not a server-side one ---
tools like Postman or \texttt{curl} never hit this wall because they aren't
bound by same-origin restrictions.

\textbf{CORS (Cross-Origin Resource Sharing)} is the mechanism the server uses
to opt certain origins back in, by sending back headers like
\texttt{Access-Control-Allow-Origin}.

\subsection{Preflight Requests}

For ``non-simple'' requests --- anything using \texttt{PUT}, \texttt{PATCH},
\texttt{DELETE}, a JSON body, or custom headers like
\texttt{Authorization} --- the browser first sends an automatic
\texttt{OPTIONS} request called a \textbf{preflight}. It asks the server
``if I were to send this real request, would you allow it?'' before sending
the actual request. The server must respond to \texttt{OPTIONS} with the
appropriate \texttt{Access-Control-Allow-*} headers, or the browser never
sends the real request at all.

\begin{diagrambox}[Preflight Flow]
\begin{verbatim}
Browser (localhost:5173)              Server (localhost:5000)
        |                                      |
        |--- OPTIONS /api/tasks -------------->|  (preflight)
        |    Origin: localhost:5173            |
        |    Access-Control-Request-Method: PUT|
        |                                      |
        |<-- Access-Control-Allow-Origin ------|
        |<-- Access-Control-Allow-Methods -----|
        |<-- Access-Control-Allow-Credentials -|
        |                                      |
        |--- PUT /api/tasks/123 -------------->|  (actual request,
        |                                      |   only sent if
        |<-- 200 OK ---------------------------|   preflight passed)
\end{verbatim}
\end{diagrambox}

\section{Configuring CORS in Express}

\begin{lstlisting}[language=JavaScript]
const cors = require("cors");

app.use(
  cors({
    origin: "http://localhost:5173", // exact frontend origin, not "*"
    credentials: true,               // allow cookies / Authorization header
  })
);
\end{lstlisting}

\begin{warnbox}[WARNING]
\texttt{origin} cannot be the wildcard \texttt{"*"} when
\texttt{credentials: true}. The CORS spec forbids combining a wildcard origin
with credentialed requests (cookies, HTTP auth) because it would let any
website read a user's authenticated response. You must list an exact origin
(or a function that validates against an allow-list of exact origins).
\end{warnbox}

\begin{lstlisting}[language=JavaScript]
// Supporting multiple known origins (e.g. staging + production)
const allowedOrigins = [
  "http://localhost:5173",
  "https://taskapp.example.com",
];

app.use(
  cors({
    origin: (incomingOrigin, callback) => {
      if (!incomingOrigin || allowedOrigins.includes(incomingOrigin)) {
        callback(null, true);
      } else {
        callback(new Error("Not allowed by CORS"));
      }
    },
    credentials: true,
  })
);
\end{lstlisting}

\section{Common CORS Errors and Fixes}

\begin{table}[h]
\centering
\begin{tabularx}{\textwidth}{X X}
\toprule
\textbf{Error in browser console} & \textbf{Fix} \\
\midrule
``No \texttt{Access-Control-Allow-Origin} header is present on the requested
resource'' &
The server never sent CORS headers at all. Add \texttt{app.use(cors(\{...\}))}
before your routes, and make sure it runs for every request (including
\texttt{OPTIONS}). \\
\addlinespace
``The value of the \texttt{Access-Control-Allow-Origin} header in the
response must not be the wildcard \texttt{*} when the request's credentials
mode is \texttt{include}'' &
Replace \texttt{origin: "*"} with the exact frontend origin string, and keep
\texttt{credentials: true} only if you actually need cookies/auth headers
sent cross-origin. \\
\addlinespace
``The request's credentials mode is \texttt{include}, but the
\texttt{Access-Control-Allow-Credentials} header ... is not \texttt{true}'' &
Add \texttt{credentials: true} inside the \texttt{cors()} config on the
server \emph{and} set \texttt{withCredentials: true} (axios) or
\texttt{credentials: "include"} (fetch) on the frontend request. \\
\addlinespace
``Response to preflight request doesn't pass access control check: It does
not have HTTP ok status'' &
Something (an auth middleware, a rate limiter) is intercepting the
\texttt{OPTIONS} request before \texttt{cors()} can respond to it. Register
\texttt{cors()} as the very first middleware, ahead of auth checks. \\
\addlinespace
``Method \texttt{PATCH} is not allowed by \texttt{Access-Control-Allow-Methods}
in preflight response'' &
Pass an explicit \texttt{methods} array to \texttt{cors()}, e.g.\
\texttt{methods: ["GET","POST","PUT","PATCH","DELETE"]}, if you rely on
methods beyond the library's defaults. \\
\bottomrule
\end{tabularx}
\end{table}

\begin{tipbox}[TIP]
In production, read the allowed origin from an environment variable
(\texttt{process.env.CLIENT\_URL}) instead of hardcoding
\texttt{localhost:5173}, so the same code works across development,
staging, and production without edits.
\end{tipbox}
